\documentclass[11pt]{article}
\usepackage[margin=1in]{geometry}
\usepackage{amsmath,amssymb,amsthm,mathtools}
\usepackage{booktabs}
\usepackage{graphicx}
\usepackage{enumitem}
\usepackage[hidelinks]{hyperref}
\hypersetup{pdftitle={Static condensation and the trace structure of ultraweak DPG for Vlasov}}

\newtheorem{proposition}{Proposition}
\newtheorem{lemma}[proposition]{Lemma}
\newtheorem{conjecture}[proposition]{Conjecture}
\theoremstyle{remark}
\newtheorem{remark}[proposition]{Remark}
\theoremstyle{definition}
\newtheorem{definition}[proposition]{Definition}

\newcommand{\bbeta}{\boldsymbol{\beta}}
\newcommand{\bn}{\boldsymbol{n}}
\newcommand{\bz}{\boldsymbol{z}}
\newcommand{\bu}{\boldsymbol{u}}
\newcommand{\bv}{\boldsymbol{v}}
\newcommand{\bE}{\boldsymbol{E}}
\newcommand{\bB}{\boldsymbol{B}}
\newcommand{\ba}{\boldsymbol{a}}
\newcommand{\bsigma}{\boldsymbol{\sigma}}
\newcommand{\that}{\hat{t}}
\newcommand{\Shat}{\widehat{S}}
\newcommand{\R}{\mathbb{R}}
\newcommand{\T}{\mathsf{T}}
\newcommand{\Tr}{\mathcal{T}}
\newcommand{\Gin}{\Gamma_{\mathrm{in}}}
\newcommand{\Gout}{\Gamma_{\mathrm{out}}}
\newcommand{\ip}[2]{\left(#1,#2\right)}
\newcommand{\dual}[2]{\left\langle #1,#2\right\rangle}
\DeclareMathOperator{\diag}{diag}
\DeclareMathOperator{\divz}{div}

\title{Static condensation and the trace structure of\texorpdfstring{\\}{} ultraweak DPG for Vlasov:\texorpdfstring{\\}{} toward a matrix-free planar MITL solver}
\author{}
\date{November 19, 2026}

\begin{document}
\maketitle

\begin{abstract}
These notes work out, as explicitly as I can, (i) how static condensation of the ultraweak DPG system for the
space--time Vlasov equation is carried out matrix-free with the separable graph norm; (ii) what the resulting
trace-only problem is, using the fact that the phase-space flow is divergence-free; and (iii) what that structure
suggests for preconditioning.  The main point is an exact algebraic identity: once each cell's faces are split into
``inflow'' and ``outflow'' sets, the condensed operator is a sum $R^\T W R + Q$ in which $R$ is a block
unit-lower-triangular \emph{sweep} operator and $Q\ge 0$ collects local inflow Schur complements that vanish in the
continuous limit.  Then $P=(R^\T W R)^{-1}$, applied by one backward and one forward sweep, satisfies
$\lambda_{\min}(P\Shat)\ge 1$, and its quality is governed by a single measurable quantity.  I also explain why the
observed $\kappa\sim h^{-2}$ growth of the current multigrid preconditioners is what one should expect from this
structure, and list experiments in order.  Statements are labelled as proved, standard, or conjectural.
\end{abstract}

\tableofcontents

%======================================================================
\section{Setting}
%======================================================================

\subsection{Equation and phase-space flow}

Phase-space--time coordinates are ordered as in the code,
\[
  \bz = (x_1,\dots,x_s,\;t,\;u_1,\dots,u_v)\in\Omega\subset\R^{n},\qquad n=s+1+v .
\]
For the planar MITL, $s=1$ and $v=2$, so $n=4$ and $\bz=(x,t,u_x,u_y)$.  For one species with charge $q$ and mass
$m$, the Vlasov equation is $\bbeta\cdot\nabla_{\bz} f = 0$ with
\[
  \bbeta(\bz) = \bigl(\bv(\bu),\;1,\;\tfrac{q}{m}\ba(\bz)\bigr),\qquad
  \ba = \bE(x,t) + \bv(\bu)\times\bB(x,t),\qquad
  \bv(\bu) = \frac{\bu}{\gamma(\bu)},\quad \gamma=\sqrt{1+|\bu|^2/c^2}.
\]
(Components of $\bv\times\bB$ not represented in the $v$ velocity dimensions are dropped, as in the planar reduction.)

\begin{lemma}[Liouville; proved]\label{lem:div}
$\divz_{\bz}\bbeta = 0$.  Moreover $\beta_t\equiv 1$, and $\beta_{x_d}=v_d(\bu)$ does not depend on $(x,t)$.
\end{lemma}
\begin{proof}
$\partial_{x_d} v_d(\bu)=0$ and $\partial_t 1 = 0$.  For the velocity part, $\bE$ does not depend on $\bu$, and
$\bv = \nabla_{\bu}\varepsilon$ with $\varepsilon(\bu)=c^2\gamma(\bu)$.  Hence
$\nabla_{\bu}\cdot(\bv\times\bB) = \bB\cdot(\nabla_{\bu}\times\nabla_{\bu}\varepsilon)
 - \nabla_{\bu}\varepsilon\cdot(\nabla_{\bu}\times\bB) = 0$, since a gradient is curl-free and $\bB$ does not
depend on $\bu$.  (The same holds after restricting to the planar components, since the restricted $\bv$ is still the
gradient of $\varepsilon$ restricted to the plane.)
\end{proof}

Note that $\partial_{u_d}\beta_{u_d}$ is \emph{not} zero individually in the relativistic case ($\partial v_y/\partial
u_x = -u_x u_y/(\gamma^3c^2)$); only the total divergence vanishes.  Nonrelativistically each $\beta_{z_d}$ is
independent of $z_d$.

\subsection{Ultraweak DPG formulation}

The mesh $\Omega_h$ consists of axis-aligned boxes $K$.  The skeleton is $\Gamma_h=\bigcup_K\partial K$.  Each face
$F$ is orthogonal to one coordinate direction $d(F)\in\{1,\dots,n\}$.

\paragraph{Unknowns.}
\begin{itemize}[nosep]
  \item Field $f\in U_f = \prod_K U_f(K)$, with $U_f(K)$ polynomials of $L^2$-degree $p-1$ (serendipity or full
        tensor), discontinuous across cells.
  \item Flux $\that\in U_t$, the discrete normal trace of $\bsigma=\bbeta f$.  In Camellia, \texttt{fluxVar} defaults
        to the space \texttt{L2} on sides, so $U_t=\prod_F U_t(F)$: each flux degree of freedom lives on exactly one
        face, and there is no coupling through face boundaries (edges, vertices of faces).  Each global face carries
        a fixed orientation; a cell sees $\that$ with sign $\epsilon_{K,F}=\pm1$ (the side parity) so that locally
        $\that$ is the flux in the direction of the outward normal $\bn_K$.
\end{itemize}
(With positivity enabled there is also the cell-constant latent field $\psi$; it is cell-local and is condensed in the
same way as $f$, but the problem is then nonlinear and I leave it aside.)

\paragraph{Test space.}  Broken: $V=\prod_K V(K)$, with $V(K)=Q_{p+\Delta p}(K)$, the full tensor-product space
(H$^1$ order $p+\Delta p$).

\paragraph{Bilinear form.}  Matching the terms in \texttt{RelativisticVlasovSpaceTimeFormulation},
\[
  b\bigl((f,\that),w\bigr) = \sum_K b_K\bigl((f,\that),w\bigr),\qquad
  b_K\bigl((f,\that),w\bigr) = -\ip{f}{\bbeta\cdot\nabla w}_K + \dual{\that}{w}_{\partial K}.
\]
Integrating by parts and using $\divz\bbeta=0$, for smooth $f$ the exact flux is $\that = (\bbeta\cdot\bn_K) f$ on
$\partial K$.

\paragraph{Test norm (separable surrogate).}  On each cell,
\[
  \|w\|_{V(K)}^2 = \alpha_K\|w\|_K^2 + \sum_{d=1}^n s_{K,d}^2\,\|\partial_d w\|_K^2 \;-\; c_K\Bigl(\int_K w\Bigr)^2,
\]
where $s_{K,d}\ge\sup_K|\beta_d|$ (so by Cauchy--Schwarz the separable norm dominates the graph norm up to the factor
$n$), $\alpha_K$ is the transport-scaled $L^2$ weight, and $c_K\ge0$ is the constant-mode weight.  Its Gram matrix
is a Kronecker sum plus a rank-one term:
\[
  G_K = G_K^{\otimes} - c_K\, \mathbf{m}\mathbf{m}^\T,\qquad
  G_K^{\otimes} = J_K\Bigl(\alpha_K \widehat M^{\otimes n} + \sum_d \xi_{K,d}\,\widehat K_{(d)}\Bigr).
\]
Here $\widehat M$, $\widehat K$ are the 1D reference mass and stiffness matrices, $\widehat K_{(d)}$ has $\widehat K$ in
slot $d$ and $\widehat M$ elsewhere, $J_K$ is the cell Jacobian, and $\xi_{K,d}=s_{K,d}^2/h_{K,d}^2$ (half-widths).

%======================================================================
\section{The DPG system in eigenbasis coordinates}
%======================================================================

Let $V_1$ solve the 1D generalized eigenproblem $\widehat K V_1 = \widehat M V_1\Lambda_1$, $V_1^\T\widehat M V_1=I$,
and let $\mathcal V = V_1^{\otimes n}$.  Then, exactly,
\[
  (G_K^{\otimes})^{-1} = \mathcal V\, D_K^{-1}\,\mathcal V^\T,\qquad
  (D_K)_{\mathbf j} = J_K\Bigl(\alpha_K + \sum_d \xi_{K,d}\,\lambda_{j_d}\Bigr).
\]
With $\mathbf z_K = (G_K^\otimes)^{-1}\mathbf m$ and $\delta_K = 1-c_K\mathbf m^\T\mathbf z_K>0$, Sherman--Morrison
gives $G_K^{-1} = (G_K^\otimes)^{-1} + (c_K/\delta_K)\,\mathbf z_K\mathbf z_K^\T$.

\begin{definition}[Eigenbasis coordinates]
For a test-space vector $\mathbf b\in\R^{n_V}$ ($n_V=\dim V(K)$), let
\[
  \Omega_K\mathbf b = \begin{bmatrix} D_K^{-1/2}\mathcal V^\T\mathbf b \\ \sqrt{c_K/\delta_K}\;\mathbf z_K^\T\mathbf b\end{bmatrix}
  \in\R^{n_V+1},
  \qquad\text{so that}\qquad \mathbf b^\T G_K^{-1}\mathbf b' = (\Omega_K\mathbf b)\cdot(\Omega_K\mathbf b').
\]
\end{definition}

Write the local matrix of $b_K$ (rows: test functions, columns: trial functions) as
$B_K = [\,B_{f}\;\;B_{t}\,]$, with
\[
  (B_f)_{ij} = -\ip{\varphi_j}{\bbeta\cdot\nabla\psi_i}_K,\qquad
  (B_t)_{ij} = \epsilon_{K,F(j)}\,\dual{\chi_j}{\psi_i}_{F(j)},
\]
and the local load $\mathbf l_K$ (from the right-hand side; zero for homogeneous Vlasov except through boundary data).
Define
\[
  W_K = \Omega_K B_K = [\,W_f\;\;W_t\,]\in\R^{(n_V+1)\times(n_f+n_t)},\qquad \mathbf e_K = \Omega_K\mathbf l_K.
\]
The DPG solution minimizes the residual
\[
  \mathcal J(f,\that) = \sum_K \bigl\| W_f\,\mathbf f_K + W_t\,\that_K - \mathbf e_K\bigr\|^2 ,
\]
where $\mathbf f_K$ and $\that_K = E_K\that$ are the cell's coefficient vectors ($E_K$ gathers, with the signs
$\epsilon_{K,F}$).  The normal equations are $S\mathbf u=\mathbf r$ with $S=\sum_K P_K^\T W_K^\T W_K P_K$; this is the
operator the current matrix-free code applies, with $W_K$ applied by sum factorization and never stored.

%======================================================================
\section{Static condensation, matrix-free}
%======================================================================

\subsection{Algebra}

For fixed $\that_K$, minimizing the cell term over $\mathbf f_K$ is a linear least-squares problem:
\[
  \mathbf f_K^\star = A_{ff}^{-1}W_f^\T(\mathbf e_K - W_t\that_K),\qquad A_{ff}=W_f^\T W_f = B_f^\T G_K^{-1}B_f.
\]
Let $\Pi_K = I - W_f A_{ff}^{-1}W_f^\T$ be the orthogonal projector onto $\operatorname{range}(W_f)^\perp$.  The
condensed functional is
\[
  \widehat{\mathcal J}(\that) = \sum_K \bigl\|\Pi_K(W_t\,\that_K - \mathbf e_K)\bigr\|^2,
\]
and the condensed system is
\begin{equation}\label{eq:condensed}
  \Shat\,\that = \hat{\mathbf r},\qquad
  \Shat = \sum_K E_K^\T\,\Shat_K\,E_K,\quad \Shat_K = W_t^\T\Pi_K W_t = A_{tt} - A_{tf}A_{ff}^{-1}A_{ft},\quad
  \hat{\mathbf r} = \sum_K E_K^\T W_t^\T\Pi_K\mathbf e_K .
\end{equation}
After solving, $\mathbf f_K$ is recovered cell by cell from $\mathbf f_K^\star$.

\subsection{Matrix-free application}

To apply $\Shat_K$ to $\that_K$:
\begin{enumerate}[nosep]
  \item $\mathbf y = W_t\that_K$: the existing flux-to-test sum-factorized path, then $\Omega_K$ (mode products with
        $V_1^\T$, a diagonal scaling, one dot product).
  \item $\mathbf g = W_f^\T\mathbf y$: the existing field path, transposed.
  \item $\mathbf s = A_{ff}^{-1}\mathbf g$: a dense Cholesky solve with a stored factor.
  \item Return $W_t^\T(\mathbf y - W_f\mathbf s)$.
\end{enumerate}
The cost is about one current cell application plus one more application of $W_f$, and an $n_f\times n_f$ triangular
solve pair.  Storage is $\tfrac12 n_f^2$ per cell for the Cholesky factor of $A_{ff}$.  With serendipity fields $n_f$
is a few hundred at most at $p=4$, so this is a fraction of the stored element-Schwarz blocks, which are
$(n_f+n_t)^2$ per cell.

\paragraph{Forming $A_{ff}$.}  $A_{ff}$ is formed once per cell by applying $W_f$ to the $n_f$ unit vectors (or, more
cheaply, by forming the dense $W_f$ once and a SYRK).  Its cost is the same kind as the current Schwarz-block setup
but over fields only.

\paragraph{Inexact alternative.}  $A_{ff}$ should be well conditioned: since $\beta_t=1$, the time derivative of the
test functions controls $f$ on every cell, so the field-only energy $f\mapsto\sup_w \ip{f}{\bbeta\cdot\nabla w}/\|w\|_V$
is expected to be equivalent to a scaled $L^2$ norm of $f$ (to be measured; not proved here).  Then $A_{ff}^{-1}$
can be replaced by a fixed-degree Chebyshev iteration preconditioned by the (diagonal, in a Legendre basis) field mass
matrix.  A fixed polynomial keeps the condensed operator linear and symmetric, so outer CG remains valid; an
adaptive inner tolerance would require flexible CG instead.

\subsection{What condensation alone does to conditioning}

\begin{proposition}[standard]\label{prop:schur}
If $S$ is symmetric positive definite and $\Shat$ is its Schur complement onto the flux block, then
$\lambda_{\min}(S)\le\lambda(\Shat)\le\lambda_{\max}(S)$; hence $\kappa(\Shat)\le\kappa(S)$.
\end{proposition}
\begin{proof}
$\Shat \le A_{tt}$ gives the upper bound; $\Shat^{-1} = (S^{-1})_{tt}$ gives $\lambda_{\max}(\Shat^{-1})\le
\lambda_{\max}(S^{-1})$.
\end{proof}
So condensation can only help, but it does not change asymptotics in $h$: a gain by a bounded factor, consistent with
experience that the gains are modest.  The value of condensation is structural, which is the subject of the rest of
these notes.

%======================================================================
\section{The trace-only problem}
%======================================================================

\subsection{Exact characterization of the local kernel (continuous broken setting)}

Consider one cell $K$ with the broken test space $H^1(K)$ (any norm on it; they are all equivalent on a fixed cell).
For $\that\in H^{-1/2}(\partial K)$ consider
\[
  \ell_{f,\that}(w) = -\ip{f}{\bbeta\cdot\nabla w}_K + \dual{\that}{w}_{\partial K}.
\]

\begin{proposition}[proved]\label{prop:kernel}
There exists $f\in L^2(K)$ with $\ell_{f,\that}=0$ in $H^1(K)'$ if and only if there is $f\in L^2(K)$ with
$\bbeta\cdot\nabla f=0$ in $\mathcal D'(K)$ and $\that = \gamma_{\bn}(\bbeta f)$, the normal trace of $\bbeta f\in
H(\mathrm{div},K)$.
\end{proposition}
\begin{proof}
Taking $w\in C_c^\infty(K)$ in $\ell_{f,\that}(w)=0$ gives $\divz(\bbeta f)=0$ in $\mathcal D'(K)$, i.e.\ $\bbeta f\in
H(\mathrm{div},K)$ with zero divergence; since $\divz\bbeta=0$, $\divz(\bbeta f)=\bbeta\cdot\nabla f$ in the
distributional sense.  Then for every $w\in H^1(K)$, Green's formula in $H(\mathrm{div},K)$ gives
$-\ip{\bbeta f}{\nabla w}_K = \ip{\divz(\bbeta f)}{w}_K - \dual{\gamma_{\bn}(\bbeta f)}{w}_{\partial K}
= -\dual{\gamma_{\bn}(\bbeta f)}{w}_{\partial K}$, so $\ell_{f,\that}(w) = \dual{\that-\gamma_{\bn}(\bbeta f)}{w}_{\partial K}$.
This vanishes for all $w\in H^1(K)$ iff $\that=\gamma_{\bn}(\bbeta f)$, because the trace map $H^1(K)\to
H^{1/2}(\partial K)$ is onto.  The converse is the same computation read backwards.
\end{proof}

So the kernel of the continuous condensed form $\that\mapsto\inf_f\|\ell_{f,\that}\|^2_{V(K)'}$ is exactly the set
of fluxes of solutions of the homogeneous transport equation in $K$.

\subsection{The kernel is a graph over the inflow data}

Let $\Gin(K)=\{\bz\in\partial K:\bbeta\cdot\bn_K<0\}$ and $\Gout(K)=\{\bbeta\cdot\bn_K>0\}$.  Because $\beta_t=1$,
every characteristic through a point of $K$ enters $K$ through $\Gin(K)$ and leaves through $\Gout(K)$ within time
$\Delta t_K$ (the cell's extent in $t$); no characteristic is trapped.  A solution of $\bbeta\cdot\nabla f=0$ is
constant along characteristics, so it is determined by its values on $\Gin$, and its outflow flux is determined by its
inflow flux:
\[
  \that|_{\Gout} = \Tr_K\bigl(\that|_{\Gin}\bigr),
\]
where $\Tr_K$ transports $f=\that/(\bbeta\cdot\bn)$ along characteristics and multiplies by $\bbeta\cdot\bn$ on the
outflow side.  The kernel is the graph of $\Tr_K$.

\begin{lemma}[standard trace theory; stated, not proved here]\label{lem:trace}
For $f$ in the graph space $\{f\in L^2(K):\bbeta\cdot\nabla f\in L^2(K)\}$, with no trapped characteristics, the trace
of $f$ lies in $L^2(\partial K,|\bbeta\cdot\bn|)$.  Equivalently, the flux $\that=(\bbeta\cdot\bn)f$ lies in the
weighted space $L^2(\partial K,|\bbeta\cdot\bn|^{-1})$, with norm $\|\that\|^2 = \int_{\partial K}|\that|^2/|\bbeta\cdot\bn|$.
(See the classical work of Bardos on traces for transport, and the analysis of DPG for transport by Broersen,
Dahmen and Stevenson.)
\end{lemma}

\begin{lemma}[proved, given smoothness]\label{lem:isometry}
$\Tr_K$ is an isometry $L^2(\Gin,|\bbeta\cdot\bn|^{-1})\to L^2(\Gout,|\bbeta\cdot\bn|^{-1})$.
\end{lemma}
\begin{proof}
If $\bbeta\cdot\nabla f=0$ then $\divz(\bbeta f^2) = f^2\divz\bbeta + 2f\,\bbeta\cdot\nabla f = 0$, so
$0=\int_{\partial K}(\bbeta\cdot\bn)f^2 = \int_{\Gout}|\bbeta\cdot\bn|f^2 - \int_{\Gin}|\bbeta\cdot\bn| f^2$, and
$|\bbeta\cdot\bn|f^2 = |\that|^2/|\bbeta\cdot\bn|$.
\end{proof}

\paragraph{Consequence for the functional setting.}  For transport the flux is not really an $H^{-1/2}$ object: the
natural trace norm is a weighted $L^2$ norm, face by face.  The condensed form is (in the continuous limit) a norm of
the distance of $\that$ from the graph of $\Tr_K$, for instance $\|\that|_{\Gout}-\Tr_K\that|_{\Gin}\|$, which by
Lemma~\ref{lem:isometry} is the same as $\|\Tr_K^{-1}\that|_{\Gout}-\that|_{\Gin}\|$.  There is no fractional-order
difficulty; all of the difficulty is directional.

\subsection{Trace classes in the Vlasov setting}

On a box mesh, the face classes are determined by the orthogonal direction.  On a face orthogonal to direction $d$,
$\bn_K=\pm\mathbf e_d$ and $\that=\pm\beta_d f$.  All classes are normal components of the one vector field
$\bsigma=\bbeta f$; equivalently, weighted traces of $f$.

\begin{center}
\begin{tabular}{@{}llp{1.45in}p{2.0in}@{}}
\toprule
face normal & flux $\that$ (outward $+\mathbf e_d$) & $\beta_d$ depends on & sign of $\beta_d$ on the face\\
\midrule
$t$ & $f$ & nothing ($\beta_t=1$) & always $+$: bottom is inflow, top is outflow \\
$x_d$ & $v_d(\bu)\,f$ & $\bu$ only (tangential) & $\operatorname{sign}(u_d)$; constant if $u_d=0$ is a mesh plane\\
$u_d$ & $\frac{q}{m}a_d(\bz)\,f$ & $x,t$, other $u$; also $u_d$ via $\gamma$ & can change within a face \\
\bottomrule
\end{tabular}
\end{center}

\begin{itemize}
  \item \textbf{$t$-faces carry the state.}  Their flux is the value of $f$ itself.  The bottom face of every cell is
        inflow and the top is outflow, without exception.  This gives a causal ordering of the cells by time layer
        inside a slab.
  \item \textbf{$x$-faces carry the spatial particle flux} $v_x f$.  By Lemma~\ref{lem:div}, $\beta_x$ is the same
        function of $\bu$ on every face of a stack of $x$-faces, so the inflow/outflow pattern is coherent along the
        stack.  If $u_x=0$ is a mesh plane, each $x$-face is entirely inflow or entirely outflow.
  \item \textbf{$u$-faces carry the acceleration flux.}  The sign of $a_d$ may change within a face, and the magnetic
        rotation $\bv\times\bB$ in the $(u_x,u_y)$ plane produces cyclic dependencies between cells of one time
        layer (Section~\ref{sec:cycles}).
\end{itemize}

%======================================================================
\section{An exact sweep factorization of the condensed operator}
%======================================================================

\subsection{Local identity}

Fix a cell $K$ and split its flux unknowns into an ``in'' set $\mathcal I_K$ and an ``out'' set $\mathcal O_K$, by
faces.  The split is a choice; ideally it matches the sign of $\bbeta\cdot\bn_K$ on each face.  Partition
\[
  \Shat_K = \begin{bmatrix}\Shat_{ii} & \Shat_{io}\\ \Shat_{oi} & \Shat_{oo}\end{bmatrix}
\]
and assume $\Shat_{oo}$ is positive definite.  (In the continuous setting a compatible flux with zero inflow is zero,
so the out block has no kernel; discretely this is to be checked.)  Define the \emph{discrete cell transport map} and
the \emph{inflow Schur complement}
\[
  \Tr_K := -\Shat_{oo}^{-1}\Shat_{oi},\qquad \Sigma_K := \Shat_{ii}-\Shat_{io}\Shat_{oo}^{-1}\Shat_{oi}\;\succeq 0 .
\]
$\Tr_K\that_i$ is exactly the outflow flux that the local DPG problem on $K$ produces from the inflow flux $\that_i$:
it minimizes $\Shat_K$ over $\that_o$ with $\that_i$ fixed, i.e.\ it is the local DPG solve with inflow data.

\begin{proposition}[proved; block $LDL^\T$]\label{prop:local}
For every $\that_K=(\that_i,\that_o)$,
\[
  \that_K^\T\Shat_K\that_K = \bigl\|\that_o-\Tr_K\that_i\bigr\|^2_{\Shat_{oo}} + \bigl\|\that_i\bigr\|^2_{\Sigma_K}.
\]
\end{proposition}
\begin{proof}
Expand: $\|\that_o-\Tr_K\that_i\|^2_{\Shat_{oo}} = \that_o^\T\Shat_{oo}\that_o + 2\that_o^\T\Shat_{oi}\that_i
+ \that_i^\T\Shat_{io}\Shat_{oo}^{-1}\Shat_{oi}\that_i$, and add $\that_i^\T\Sigma_K\that_i$.
\end{proof}

\begin{remark}
In the continuous limit with the ideal split, Proposition~\ref{prop:kernel} and the graph structure give
$\Sigma_K=0$: for any inflow data, the compatible outflow data make the residual zero.  Discretely $\Sigma_K\ne0$ for
two reasons: the polynomial spaces do not contain the exact transported solutions, and the chosen split may disagree
with the sign of $\bbeta\cdot\bn$ on some faces.  The size of $\Sigma_K$ relative to the first term is the quantity
that decides everything below.
\end{remark}

\subsection{Global assembly: ownership and the sweep operator}

Require that every interior face be ``out'' for exactly one of its two cells and ``in'' for the other.  Call the first
cell its \emph{owner}.  Boundary faces whose fluxes are prescribed (Dirichlet data, in particular all inflow
boundary faces) carry data; every other boundary face is owned by its only cell.  Then
every flux unknown has exactly one owner, and the unknowns are partitioned by owner:
$\that = (\that_{\mathcal O_K})_K$.

\begin{definition}
Let $R$ be the block matrix with block rows and columns indexed by cells, defined by
\[
  (R\that)_K = \that_{\mathcal O_K} - \Tr_K\,\that_{\mathcal I_K},
\]
where $\that_{\mathcal I_K}$ consists of unknowns owned by $K$'s upwind neighbours (boundary data contribute to the
right-hand side).  Let $W=\diag_K(\Shat_{oo,K})$ and $Q=\sum_K E_{\mathcal I_K}^\T\Sigma_K E_{\mathcal I_K}$.
\end{definition}

\begin{proposition}[proved]\label{prop:global}
$\Shat = R^\T W R + Q$, with $Q\succeq0$.  The diagonal blocks of $R$ are identities.  If the directed graph ``$K'\to
K$ when $K'$ owns a face in $\mathcal I_K$'' is acyclic, then in a topological order $R$ is block unit
lower triangular, and $R^{-1}$ is applied by a forward sweep:
\[
  \text{for $K$ in upwind-to-downwind order:}\qquad \that_{\mathcal O_K} = \mathbf g_K + \Tr_K\,\that_{\mathcal I_K}.
\]
\end{proposition}
\begin{proof}
Sum Proposition~\ref{prop:local} over cells; each unknown appears as an out unknown exactly once.  The rest is
forward substitution.
\end{proof}

\subsection{The sweep preconditioner}

Let $P=(R^\T WR)^{-1} = R^{-1}W^{-1}R^{-\T}$.  Applying $P$ to a residual $\mathbf r$:
\begin{enumerate}[nosep]
  \item Backward sweep (downwind to upwind): solve $R^\T\mathbf y=\mathbf r$, i.e.\ for $K$ in reverse order,
        $\mathbf y_K = \mathbf r_K + \sum_{K'' :\,K\text{ owns a face in }\mathcal I_{K''}}
        \bigl(\Tr_{K''}^\T\mathbf y_{K''}\bigr)\big|_{\mathcal O_K}$.
  \item Cell-local solves: $\mathbf y_K\leftarrow\Shat_{oo,K}^{-1}\mathbf y_K$.
  \item Forward sweep (upwind to downwind): solve $R\mathbf x=\mathbf y$ as in Proposition~\ref{prop:global}.
\end{enumerate}
$P$ is symmetric positive definite and fixed, so it can be used in CG or as a multigrid smoother.

\begin{proposition}[proved]\label{prop:bound}
$\lambda_{\min}(P\Shat)\ge1$ and
\[
  \kappa(P\Shat)\le\lambda_{\max}(P\Shat) = 1+\rho,\qquad
  \rho = \sup_{\that\ne0}\frac{\that^\T Q\,\that}{\that^\T R^\T WR\,\that}.
\]
\end{proposition}
\begin{proof}
$\Shat = R^\T WR + Q\succeq R^\T WR$, and the Rayleigh quotient of $P\Shat$ in the $R^\T WR$ inner product is
$1+\that^\T Q\that/\that^\T R^\T WR\that$.
\end{proof}

\begin{conjecture}\label{conj:rho}
With the split chosen by the sign of $\bbeta\cdot\bn$ (face-averaged where it changes sign), $\rho$ is bounded
independently of $h$, and is small when the cell Courant numbers $\nu_{K,d}=|\beta_d|\Delta t_K/h_{K,d}$ are small
on the faces where the split disagrees with the flow.
\end{conjecture}

This is the claim to test.  If it holds, sweeps give $h$-independent iteration counts, which no local (additive)
smoother can, and $\rho$ is a single computable number per mesh: the largest eigenvalue of a generalized eigenproblem
$Q\mathbf x=\mu R^\T WR\mathbf x$, estimated by Lanczos within PCG exactly as we already estimate $\kappa$.

\subsection{Equivalent uncondensed form}

Because fields are cell-local, the same preconditioner can be written without condensing: for each cell, given the
inflow fluxes, solve the local DPG problem for $(f_K,\that_{\mathcal O_K})$.  This is a sweep of local DPG solves,
the DPG analogue of an upwind DG sweep.  The condensed form is cleaner to analyse and cheaper to store.

\subsection{Cost}

Per cell: a Cholesky factor of $\Shat_{oo}$ ($\tfrac12 n_o^2$) and $\Tr_K$ ($n_o\times n_i$), with $n_i+n_o=n_t$.  One
application of $P$ costs two products with $\Tr_K$ or $\Tr_K^\T$ and one triangular solve pair per cell, i.e.\
$O(n_t^2)$ flops per cell.  That is far cheaper than one matrix-free application of $\Shat$.  The blocks are formed
once from $\Shat_K$, which is assembled densely per cell from the eigenbasis coordinates
($\Shat_K=W_t^\T\Pi_KW_t$); transient memory is one $(n_V+1)\times(n_f+n_t)$ matrix.  Storage is comparable to the
current element-Schwarz blocks, which measured about 2.6~MB per cell at $p=4$.

%======================================================================
\section{Cycles, mixed faces, and time-step size}\label{sec:cycles}
%======================================================================

\paragraph{Time layers are always acyclic.}  The bottom face is in and the top face is out for every cell, so the
dependency graph restricted to $t$-faces goes strictly upward.  Cycles can only come from lateral faces within one
time layer.

\paragraph{Magnetic rotation.}  With $\bE=0$ and $\bB=B\mathbf e_z$, $\beta_{(u_x,u_y)}=\frac{qB}{m}(v_y,-v_x)$:
circular flow about $\bu=0$.  Four cells of one time layer around the origin in $(u_x,u_y)$ then feed one another in a
cycle.  The same happens elsewhere in the gyration region.

\paragraph{Breaking cycles costs only $\rho$.}  Proposition~\ref{prop:global} needs only an acyclic ownership; the
split need not match the flow.  One can always choose an acyclic split, e.g.\ order cells within a layer and give each
lateral face to the earlier cell.  A face assigned against the flow does not break the identity
$\Shat=R^\T WR+Q$; it increases $\Sigma_K$ and hence $\rho$.  Mixed-sign faces are treated the same way.

\paragraph{Heuristic size of the penalty.}  A characteristic entering through the bottom of $K$ crosses a lateral
face orthogonal to $d$ only if it travels farther than the distance to that face in time $\Delta t_K$.  The fraction of
the flux that leaves $K$ laterally is therefore about $\min\bigl(1,\sum_d\nu_{K,d}\bigr)$.  A cycle through $k$ cells
passes on a fraction of order $\nu^k$.  For thin time cells ($\nu\ll1$) the lateral coupling, and hence the penalty for
misassigned faces, is small.  For magnetized electrons in the $u$ directions,
$\nu_{K,u}\approx\omega_c\Delta t_K\,|\bu|/h_{K,u}$ (nonrelativistically), i.e.\ the gyro-angle per time cell times
$|\bu|/h_u$.  This may not be small in the MITL; it should be computed for the meshes we use.

\paragraph{Parallelism.}  Sweeps are sequential within a layer.  Standard remedies from neutron transport (KBA
pipelining across subdomains, lagging across processor boundaries) apply.  Lagging across subdomains is again a
choice of split, so it shows up as a contribution to $\rho$.

%======================================================================
\section{\texorpdfstring{Why the current preconditioners show $\kappa\sim h^{-2}$}{Why the current preconditioners show kappa ~ h**(-2)}}
%======================================================================

Proposition~\ref{prop:global} says that $\Shat$ is, up to $Q$, a weighted least-squares form of the sweep equations
$R\that=\mathbf g$.  Least-squares forms of a first-order transport operator behave like a second-order operator
along $\bbeta$ that is nearly decoupled across $\bbeta$.  Their low-energy modes are smooth along characteristics but
may be rough across them.  Isotropic $h$-coarsening cannot represent such modes, and additive local smoothers (point,
trace-block, element Schwarz, Chebyshev-accelerated or not) cannot remove them.  The expected result is the behaviour
of a method without coarse correction in the direction along $\bbeta$, i.e.\ $\kappa\sim h^{-2}$.  That is what we
measure, for matrix-free and production GMG alike:

\begin{center}
\begin{tabular}{@{}lrrlll@{}}
\toprule
case & cells & dofs & MF element Schwarz & MF trace-block & production GMG\\
 & & & its ($\kappa$) & its ($\kappa$) & its\\
\midrule
$p=2$, $r=1$ & 128  & 5\,344   & 9 (1.5)  & 19 (4.3)  & 47\\
$p=2$, $r=2$ & 2048 & 76\,544  & 20 (4.9) & 39 (16.9) & 88\\
$p=4$, $r=1$ & 128  & 17\,216  & 11 (1.9) & 24 (6.4)  & 71\\
$p=4$, $r=2$ & 2048 & 247\,296 & 23 (6.5) & 47 (25.2) & (too large here)\\
\bottomrule
\end{tabular}
\end{center}

Base mesh $2\times1\times2\times2$ in $(x,t,u_x,u_y)$, $r$ uniform refinements, $\Delta p=2$, separable norm, full GMG
$h$-hierarchy to the coarsest mesh at $p=1$, Chebyshev degree 4 for the matrix-free runs.  Each refinement multiplies
$\kappa$ by about 3.5--4 and doubles the iterations.

This explanation is a hypothesis.  It predicts that the eigenvector of the smallest eigenvalue of the preconditioned
operator is smooth along $\bbeta$ and oscillatory across it, which is directly checkable (Experiment~1 below).

%======================================================================
\section{Experiments, in order}
%======================================================================

\begin{enumerate}
  \item \textbf{Diagnose the slow mode.}  From the PCG Lanczos data, recover the Ritz vector for $\lambda_{\min}$ of
        the current GMG-preconditioned operator, and measure its variation along and across $\bbeta$ (e.g.\ the
        ratio $\|\bbeta\cdot\nabla f\|/\|\nabla_\perp f\|$ of the recovered field).  Cheap; tests the hypothesis of
        the previous section.
  \item \textbf{Sweeps without condensation.}  Symmetric multiplicative element Schwarz, ordered by time layer and
        then by an acyclic upwind order within the layer, using the element blocks we already store.  Matrix-free
        residual updates: after correcting cell $K$, only $K$ and its face neighbours' contributions to the residual
        change, so one keeps per-cell eigenbasis coordinates $\mathbf y_{K'}=W_{K'}\mathbf u_{K'}$ and updates them on
        shared faces.  Use it as a preconditioner alone and as the GMG smoother; record $\kappa$ for $r=1,2,3$.
  \item \textbf{Condensation and the sweep preconditioner $P$.}  Implement (\ref{eq:condensed}) matrix-free; form
        $\Tr_K$, $\Shat_{oo,K}$, $\Sigma_K$ per cell; apply $P$ by two sweeps.  Measure $\rho$
        (Proposition~\ref{prop:bound}) and $\kappa(P\Shat)$ under refinement, and their dependence on $\Delta t_K$ and
        on the cell Courant numbers.  Compare face-block Jacobi on $\Shat$ as the local baseline.
  \item \textbf{Combine with coarse correction.}  If $\rho$ grows (because of cycles or under-resolution), use $P$ as
        the smoother in a trace-only GMG hierarchy, so that the coarse space handles what the sweep leaves.
        \texttt{GMGOperator} already supports static condensation.
  \item \textbf{Production.}  If Experiments 2 and 3 show bounded $\kappa$, an ordered multiplicative Schwarz
        smoother
        for production GMG is a small change relative to the existing additive Schwarz, and would test the same
        idea at scale.
\end{enumerate}

%======================================================================
\section{First measurements (Experiments 2 and 3, uncondensed)}\label{sec:measurements}
%======================================================================

\paragraph{What was run.}  All runs: $p=2$, $\Delta p=2$, separable norm, full-tensor test space, serendipity trial,
one core, PCG to relative residual $10^{-10}$, \emph{one level} (no coarse space).  The operator is that of the first
time slab before its solve.  In that slab the fields are zero, so $\beta_{u_x}=\beta_{u_y}=0$: transport is in $x$ and
$t$ only, the $u$-faces are characteristic ($\bbeta\cdot\bn=0$), and the upwind graph has no cycles.  Entries are
``iterations ($\kappa$)''.  Code: \texttt{ElementMultiplicativeSchwarz}, \texttt{upwindCellOrder()} and
\texttt{transportSweep()} in \texttt{MatrixFreeDPGOperator}; driver option \texttt{--matrixFreeMultiplicativeSchwarz}.

\begin{itemize}[nosep]
  \item \emph{Element blocks}: each cell's fields and all its face fluxes (overlapping), $A_K=R_KSR_K^\T$.
  \item \emph{Owner blocks}: each cell's fields and the fluxes it owns (non-overlapping), $A_K=R_KSR_K^\T$.
  \item \emph{Multiplicative}: symmetric block Gauss--Seidel (forward then backward sweep) in the given order.
  \item \emph{Transport sweep}: $P=R^{-1}W^{-1}R^{-\T}$ of Proposition~\ref{prop:global}, uncondensed: $W$ and $\Tr_K$
        come from each cell's \emph{own} $S_K$ split into owned (fields, outflow fluxes) and inflow parts.
\end{itemize}

\paragraph{Uniform refinement.}  Base mesh $2\times1\times2\times2$ in $(x,t,u_x,u_y)$, $r$ uniform refinements.
\begin{center}
\resizebox{\textwidth}{!}{%
\begin{tabular}{@{}lrrrr@{}}
\toprule
preconditioner & $r=0$ (8 cells) & $r=1$ (128) & $r=2$ (2048) & $\kappa$ growth per $r$\\
\midrule
element blocks, additive & 19 (6.0) & 56 (33) & 117 (141) & $\times4.3$\\
element blocks, multiplicative, upwind order & 9 (1.6) & 20 (5.0) & 39 (16.3) & $\times3.2$\\
element blocks, multiplicative, cell-ID order & 9 (1.6) & 21 (5.0) & 40 (16.5) & \\
element blocks, multiplicative, random order & 9 (1.6) & 23 (5.7) & 42 (18.4) & \\
owner blocks, multiplicative, upwind order & 13 (3.1) & 35 (12.8) & 69 (48) & $\times3.8$\\
\textbf{transport sweep} $P$ & 15 (2.9) & 21 (4.5) & \textbf{24 (7.1)} & $\times1.6$\\
\midrule
for reference: matrix-free GMG, element Schwarz & --- & 9 (1.5) & 20 (4.9) & $\times3.3$\\
\bottomrule
\end{tabular}}
\end{center}
At 2048 cells the sweep also had the smallest cost of the one-level methods: 17~s solve and 21~s setup, against 31~s
and 41~s for multiplicative element blocks.

\paragraph{One direction at a time.}  From $4\times2\times4\times4$ (128 cells; ``$v$'' refines both $u$ directions).
The spatial mesh is geometrically graded toward the cathode ($x=0$); \texttt{--numSpatialRefinements} bisects every
$x$-element and keeps the grading.
\begin{center}
\resizebox{\textwidth}{!}{%
\begin{tabular}{@{}lrrr@{}}
\toprule
mesh change & additive & multiplicative (upwind) & transport sweep\\
\midrule
none (128 cells) & 56 (32) & 20 (4.9) & 22 (5.2)\\
$x$ bisected once / twice & 83 (87) / 130 (284) & 29 (10.7) / 46 (31.6) & 26 (7.4) / 32 (13.3)\\
$t$ doubled / quadrupled & 80 (63) / 120 (139) & 27 (7.9) / 38 (15.1) & 25 (6.8) / 28 (9.0)\\
$u_x,u_y$ doubled & 61 (36) & 21 (4.9) & 22 (5.8)\\
\midrule
$x$ grading deepened: 8, 16 elements ($h_{\min}=2^{-7}$, $2^{-15}$) & 93 (192), 113 (683) & 36 (24), 45 (81) & 44 (27), 236 (6036)\\
\bottomrule
\end{tabular}}
\end{center}
With 16 graded $x$-elements, shortening the slab to a $t$ width of $1$, $\tfrac12$, $\tfrac14$ and $\tfrac18$ gives
sweep condition numbers of about 6036, 3790, 2063 and 1059.  Holding the cell Courant number fixed (refining $x$ grading and $t$ together)
does not help.

\paragraph{Where the sweep's bad mode lives.}  Power iteration on $P S$ gives its top eigenvector.  On the mesh with
$x$ bisected twice ($\kappa=13.3$): 97\% of $|x|^2$ is on $u_y$-face fluxes (1.3\% $u_x$-face, 1\% each $t$- and
$x$-face, $<10^{-4}$ fields), in the bottom time layer, at the largest $u_x$ (the fastest transport in $+x$), spread
along $x$.  The same holds on the deep graded mesh, concentrated in the thin cells at the cathode.

\paragraph{What this says.}
\begin{enumerate}[nosep]
  \item The identity $S=R^\T WR+Q$ holds: $\lambda_{\min}(PS)=1$ to five digits in every run.
  \item $\rho$ is not bounded, so Conjecture~\ref{conj:rho} is false as stated.  But it grows far more slowly than for
        any local or Gauss--Seidel method (about $\times1.6$ per uniform refinement against $\times3$--$4$), and one level
        of the sweep is already competitive with the full GMG hierarchy at 2048 cells.
  \item Ordering barely matters for block Gauss--Seidel on $S$, as expected: Gauss--Seidel on the normal-equations
        operator is not the sweep.  The sweep needs $R$, i.e.\ each cell's \emph{own} local solve with inflow data.
  \item The growth is on \textbf{characteristic faces}, where $\bbeta\cdot\bn=0$ and the in/out split is undefined.
        In the continuous problem the compatible flux there is zero, and \emph{neither} neighbour's data determines it.
        The sweep must nevertheless give the face to one cell, which turns it into inconsistent inflow data for the
        other and makes $\Sigma_K$ large.  This is exactly a trace class in the sense of Section~4.3: in a field-free
        slab, the $u$-faces.
  \item Thin cells (the deep cathode grading) make the characteristic-face problem much worse.
\end{enumerate}

\paragraph{Next steps suggested by this.}
\begin{itemize}[nosep]
  \item Take characteristic-face fluxes ($|\bbeta\cdot\bn|$ small on the face) out of the sweep: sweep the other
        unknowns and treat the characteristic-face class by its own face blocks, combined multiplicatively and
        symmetrically with $P$.  Then measure $\rho$ again.
  \item Repeat on a later slab with nonzero fields, where $u$-faces carry real transport and the $\bv\times\bB$ cycles
        appear.
  \item Use the sweep as the smoother in the existing hierarchy.
\end{itemize}

\section*{Summary of what is proved and what is not}
\begin{itemize}[nosep]
  \item Proved here: Lemma~\ref{lem:div}, Propositions~\ref{prop:schur}, \ref{prop:kernel}, \ref{prop:local},
        \ref{prop:global}, \ref{prop:bound}, Lemma~\ref{lem:isometry} (for smooth solutions).
  \item Standard, cited: the weighted-$L^2$ trace space (Lemma~\ref{lem:trace}).
  \item Expected, to be measured: conditioning of $A_{ff}$; positive definiteness of $\Shat_{oo}$.
  \item Conjectural: the explanation of the observed $h^{-2}$ growth.  Conjecture~\ref{conj:rho} is contradicted by the
        measurements of Section~\ref{sec:measurements} (Experiment 3 there, uncondensed).
\end{itemize}

\end{document}
